\section{Weekly Summary}

\begin{tcolorbox}[colback=yellow!10,colframe=black!50,title=AI-Generated Content]
\noindent The summary below was written by Claude (Anthropic) from that week's regression outputs and series descriptions. It has not been reviewed or fact-checked by a human, and should be read as an automated, informal recap -- not analysis.
\end{tcolorbox}

Welcome back to the weekly ritual where we take two economic time series that have never met, have no business being in the same room, and force them into a linear regression together to see what happens. As always, the house motto applies: these pairs are chosen at random, no causal claim is being made, and spurious correlation isn't a bug here, it's basically the entire personality of the project. Please do not reallocate your retirement savings based on anything below.

Monday kicked things off with the Balance on Services (discontinued, naturally -- all our best friends are discontinued) versus the exchange rate for Belize. The raw regression coughed up a p-value of 0.045, which is the statistical equivalent of squeaking through a door right as it closes. But here's the plot twist: the detrended version actually got stronger, flipping the sign to negative and tightening up to a p-value of 0.018. So after we stripped out the shared drift, these two STILL want to hang out, except now in opposite directions. I have no theory for why the US services balance would have anything to do with the Belizean dollar per US dollar, and I refuse to invent one, but credit where it's due: this little guy survived detrending, which is more than most contestants can say.

Then came Tuesday, and oh, it was a showstopper. Gasoline prices in the South Census Region versus 30-year FHA mortgage rates. The raw regression was significant, sure, but the detrended version is where things got genuinely spooky: an R-squared of 0.63 and a p-value with fifty-one zeros after the decimal point. Fifty-one! That's not a correlation, that's a marriage. And here's the uncomfortable part -- this one actually makes a smidge of sense, because gas prices and mortgage rates both ride the same 1970s-to-1990s inflation rollercoaster, and that linkage held up even after we yanked out the trend. So this is the rare entry I'm legally obligated to call "actually kind of interesting," which pains me to type. Enjoy it, because it won't last.

Wednesday brought machinery corporations' cash piles versus government-sponsored enterprise debt, which in raw form posted a jaw-dropping R-squared of 0.80 -- before immediately face-planting into irrelevance the moment we detrended it (p-value crashing to a limp 0.13). This is the textbook trend mirage: two things that both went up over time, holding hands only because the calendar was dragging them both forward. Classic. Thursday closed us out with Sweden's government consumption prices versus Kansas City energy costs, which looked respectable raw (p = 0.0004) and then evaporated completely once detrended, landing at a gloriously meaningless p-value of 0.67. Two for two on the mirages this week, which is honestly the natural order of things. See you next week, when we once again mistake coincidence for destiny.
